\documentclass[openany,fontset=fandol]{ctexbook}
\usepackage{amsmath,amssymb}
\usepackage[unicode,hidelinks]{hyperref}
\newcommand{\note}[1]{\footnote{译注：#1}}
\begin{document}
\title{代数几何基础第五卷预备笔记\\注13.4(e)：拓扑结构}
\author{据用户上传英译编辑稿译出}
\date{2026年11月17日}
\maketitle
\chapter*{范围与校订状态}
本批连续译出底本第53页注13.4(e)全文，至“with group $\pi_1$”止。下一处为同页§14标题“Conic Projections”，随后接第54页正文；本批尚未译入§14。

根据13.1、13.2及本句后文，统一多重初等概形和关联解析空间的记号；同伦群指标按正次数解释。原写法及校订理由见脚注。普遍覆盖空间与$\mathbb C^n$同胚的原断言完整保留，但不将同胚加强成双全纯，也不将本项未给出的拓扑论证当作已经核实的证明。全部旧待核保持，未合并主稿。
\chapter*{§13 初等态射与M. Artin定理（续）}
\noindent\textbf{注13.4（续）。}

\noindent e) 13.2之所以值得关注，原因之一在于多重初等代数概形$U$的拓扑结构特别简单。\note{底本写“the elementary algebraic schemes $U$”。13.2的邻域为多重初等概形，13.1也以结构态射为初等态射之复合来定义多重初等概形；本项随后讨论一般维数$n$以及自由群的逐次扩张。因此此处正文用“多重初等”，不把$U$误限为结构态射本身是初等态射的相对一维情形。}
例如，若基域为复数域，并以$U^{\mathrm{en}}$表示与$U$关联的解析空间，则同伦群
\[
\pi_i(U^{\mathrm{en}})=0\qquad(i\geq2),
\]
而$\pi_1$是自由群的逐次扩张。\note{底本先写“if $Y^{\mathrm{en}}$ denotes the analytic space associated to $U$”，后文却始终使用$U^{\mathrm{en}}$，故将前一个$Y$订为$U$，上标“en”照存。底本同伦群指标写“$i\ne1$”；按同伦群的正次数范围，这就是$i\geq2$。$\pi_0$通常是连通分支集合而非这里的同伦群；正文不将其当作群来写等于零。未增写或指定原文未给出的扩张列。}

因此，$U^{\mathrm{en}}$是一个为$\pi_1$分类的“$K(\pi_1,1)$空间”；更确切地说，它的普遍覆盖空间同胚于$\mathbb C^n$，从而当然是可缩的，而这一覆盖是一个群为$\pi_1$的“普遍主纤维化”。\note{底本写“homeomorphic to $C^n$”，这里$C$按复数域排为$\mathbb C$，$n$沿用本节维数记号。“同胚”仅是拓扑意义，原文没有声称与$\mathbb C^n$双全纯。普遍覆盖可缩比“同胚于$\mathbb C^n$”弱，不能由前者反推后者；本项未附这项较强断言的证明。本译文保存原断言，不宣称已经独立验证其拓扑论证。末句的覆盖及群作用是作者所谓“普遍主纤维化”的语境，不另行推定一个未给出的代数主丛构造。}
\end{document}
